
The approach used for the motion control of the vehicle exploits its
holonomic design by
decoupling the translational and rotational modes. To do so, we first
apply feedback linearisation~\cite{sastry99} to the translational part
of~\eqref{eq:dynsys}:
\begin{equation}
  \label{eq:fctrl}
    \left\{
    \begin{split}
      \dot{\bar{x}} &= \bar{v} \\
      \dot{\bar{v}} &= \bar{p} \\
      \bar{F} &= m \mathbf{R}^T \bar{p}
    \end{split}
  \right.
\end{equation}
and then design a feedback controller for $\bar{p}$. Since this
dynamical system is second order, a PD controller is enough to ensure
convergence:
\begin{equation}
  \label{eq:pctrl}
  \begin{split}
    \bar{p} &= -k_x \bar{e}_x -k_v \bar{e}_v \\
    \bar{e}_x &= \bar{x} - \bar{x}_d \\
    \bar{e}_v &= \bar{v} - \bar{v}_d \\
  \end{split}
\end{equation}
where $\bar{x}_d$ and $\bar{v}_d$ are the desired position and
velocity vectors in the inertial frame~$\mathcal{I}$, and $k_x$ and
$k_v$ are the proportional and derivative gains of the PD controller.

For the attitude control we follow the convergent $SO(3)$ controller
proposed in~\cite{lee10} for the attitude of a quadrotor, negleting
the Coriolis and the angular velocity reference in the expression for
the torque, as suggested in~\cite{mahony12}:
\begin{equation}
  \label{eq:mctrl}
  \begin{split}
    \bar{M} &= -k_R\,\bar{e}_R - k_\omega\,\bar{e}_\omega \\
    \bar{e}_R &= \frac{1}{2} S^{-1}( \mathbf{R}^T_d \mathbf{R} -
    \mathbf{R}^T \mathbf{R}_d ) \\
    \bar{e}_\omega &= \omega - \mathbf{R}^T \mathbf{R}_d\,\omega_d
  \end{split}
\end{equation}
where $\mathbf{R}_d$ is the rotation matrix of the desired attitude,
$\omega_d$ is the desired angular velocity vector, with $k_R$ and
$k_\omega$ as controller gains. The $S^{-1}$ function performs the inverse
operation as the one defined in~\eqref{eq:S}, that is, recovers
the vector from a given skew-symmetric matrix.


%%% Local Variables:
%%% mode: latex
%%% TeX-master: "main"
%%% End:
